  \begin{vidu}
  	Một chất điểm dao động điều hòa với phương trình $x = 6\cos\left(2\pi t - \dfrac{\pi}{3}\right)$ (cm, s). 	
  	\begin{itemize}
  		\item[1.] Tính quãng đường vật đi được trong $5$ giây kể từ thời điểm ban đầu.
  		\item[2.] Tính quãng đường vật đi được trong $2,5$ giây kể từ thời điểm $t = 0$.
  		\item[3.] Tính quãng đường vật đi được trong $\dfrac{7}{12}$ giây kể từ thời điểm ban đầu.
  		\item[4.] Tính quãng đường vật đi được từ thời điểm $t_1 = 1,25\ s$ đến $t_2 = 4,5\ s$.
  		\item[5.] Tính tốc độ trung bình của vật trong $1,25$ giây đầu tiên.
  	\end{itemize}
  	\loigiai{
  		\begin{itemize}
  			\item[a)] 120 cm
  			\item[b)] 60 cm
  			\item[c)] $7+3\sqrt{3}$ cm
  			\item[d)] $75+3\sqrt{3}$ cm
  			\item[e)] $22,24$ cm/s
  		\end{itemize}		
  	}
  \end{vidu}
  %\dotlineEX{24}
  %\loai{Tính quãng đường vật đi được trong quãng thời gian $\Delta t=(2n+1)\dfrac{T}{2}$}
  \begin{vidu}%[3P1-13.1K][Loai1]  
  	Một chất điểm dao động điều hòa theo phương trình $x=4\sin \left({\pi t+\dfrac{\pi }{6}}\right)$ cm, t tính bằng s. Quãng đường chất điểm đi được sau 7 giây kể từ $t=0$ là
  	\choice
  	{\True 56 cm}
  	{48 cm}
  	{58 cm}
  	{54 cm}
  	\loigiai{
  		\begin{itemize}
  			\item Chu kì dao động của vật $T=\dfrac{2\pi}{\omega}=2$ s.
  			\item Ta có $\Delta t=3,5T\Rightarrow $ quãng đường vật đi được là $S=3.4A+2A=56\ cm.$
  			
  		\end{itemize}
  	}
  	%\haidong{9}
  \end{vidu} %\dotlineEX{6}
  %\loai{Quãng đường vật đi được sau thời gian $\Delta t$ kể từ thời điểm ban đầu -  tách thời gian}
  \begin{vidu}%[3P1-13.1K][Loai1]
  	Cho một chất điểm dao động điều hòa với tần số bằng 4 Hz và biên độ bằng 4 cm. Thời điểm ban đầu $t=0$ chất điểm đi qua vị trí li độ 2 cm theo chiều âm. Quãng đường đi sau 0,6875 s chuyển động của chất điểm là
  	\choice
  	{$46+2\sqrt{3}$ cm}
  	{$46-2\sqrt{3}$ cm}
  	{\True $42+2\sqrt{3}$ cm}
  	{$42 - 2\sqrt{3}$ cm}
  	\loigiai{
  		\immini{
  			\begin{itemize}
  				\item 
  				Biểu diễn vị trí ban đầu của vật bằng điểm $M_0$ trên đường tròn lượng giác như hình vẽ.
  				\item Ta có: $T = \dfrac{1}{f} = 0,25\ s \Rightarrow \omega = \dfrac{2\pi}{0,25} = 8\pi$ rad/s.
  				\item Ta lại có: $t = 0,6875\ s = 2.0,25 + 0,1875 = 2T + 0,1875$.
  				\item Sau 2T, chất điểm quay được 2 vòng và lại trở về đúng vị trí $M_0$ trên đường tròn $\Rightarrow S_1 = 2.4.4 = 32\ cm$.
  				\item Trong khoảng thời gian 0,1875 s tiếp theo, chất điểm quay được 1 góc $\Delta \varphi = \dfrac{3\pi}{2}$ và đến vị trí điểm $M_1$ trên đường tròn lượng giác.\\
  				$\Rightarrow S_2 = 2+8+2\sqrt{3} = 10+2\sqrt{3}\ cm$.
  			\end{itemize}
  		}{\begin{tikzpicture}[scale=1,thick,>=stealth']
  				\tkzInit[ymin=-3,ymax=3,xmin=-3,xmax=3]\tkzClip
  				\tkzDefPoints{-2.5/0/m, 2.5/0/n, 0/0/o, 0/2.5/p, 0/-2.5/q}
  				\tkzDefShiftPoint[o](60:2){0}
  				\tkzDefShiftPoint[o](-30:2){m'}
  				\tkzDefPointBy[projection=onto m--n](0) \tkzGetPoint{h}
  				\tkzDefPointBy[projection=onto m--n](m') \tkzGetPoint{h'}
  				\draw(o)circle(2cm);
  				\tkzDrawSegments[dashed](m',h' 0,h)
  				\tkzDrawSegments(p,q)
  				\tkzDrawSegments[->](m,n)
  				\tkzDrawSegments[blue,->](o,0 o,m')
  				\tkzDrawPoints[size=3](o,h,0,m',h')
  				\tkzMarkAngles[size=0.7cm,arrows=->](0,o,m')
  				\node at (0)[above right]{$M_0$};
  				\node at (m')[below right]{$M_1$};
  				\node at (2,0)[below right]{$4$};
  				\node at (h)[below]{$2$};
  				\node at (h')[above]{$2\sqrt{3}$};
  		\end{tikzpicture}}
  		\begin{itemize}
  			\item Vậy quãng đường chất điểm đi được sau 0,6875 s chuyển động là: 
  			\begin{align*}
  				S = S_1+ S_2 = 32+10+2\sqrt{3} = 42+2\sqrt{3}\ cm.
  			\end{align*}
  		\end{itemize}
  	}
  	%\haidong{8}
  \end{vidu} %\dotlineEX{6}
  %\loai{Tìm thời gian vật đi được quãng đường $S < 2A$}
  \begin{vidu}%[3P1-13.2K][Loai1] 
  	Vật dao động điều hòa theo phương trình: $x = \cos(\pi t - \dfrac{2\pi}{3})$ dm. Thời gian vật đi được quãng đường $S = 5$ cm kể từ thời điểm ban đầu ($t = 0$) là
  	\choice
  	{ $\dfrac{1}{9}$ s}
  	{$\dfrac{1}{3}$ s}
  	{\True $\dfrac{1}{6}$ s}
  	{$\dfrac{7}{3}$ s}
  	\loigiai{
  		\begin{itemize}
  			\item $x = 10\cos(\pi+t-\dfrac{2\pi}{3})\ cm$.
  			\item Lúc đầu vật ở li độ $x = -\dfrac{A}{2}$ và đang đi theo chiều dương.
  			\item Vật đi được quãng đường $s = \dfrac{A}{2}\Rightarrow t = \dfrac{T}{12} = \dfrac{T}{12} = \dfrac{1}{6}\ s$.
  			
  		\end{itemize}
  	}
  	%\haidong{6}
  \end{vidu} %\dotlineEX{6}
  %\loai{Tìm thời gian vật đi được quãng đường $2A<S < 4A$}
  \begin{vidu}%[3P1-13.2K][Loai1] 
  	Một chất điểm dao động điều hòa với chu kì 3 s trên một quỹ đạo dài 4 cm. Biết thời điểm ban đầu vật ở vị trí cân bằng theo chiều dương, thời gian để chất điểm đi được 5 cm kể từ thời điểm ban đầu là
  	\choice
  	{$\dfrac{15}{12}$ s}
  	{2 s}
  	{\True $\dfrac{7}{4}$ s}
  	{$\dfrac{18}{12}$ s}
  	\loigiai{
  		\immini{\begin{itemize}
  				\item 
  				Con lắc lò xo dao động điều hòa trên quỹ đạo dài 4 cm nên biên độ của nó là 2 cm.
  				\item Ta có $5 = 2.2 + 1 = 2A + 1\Rightarrow t = \dfrac{T}{2} + \Delta t$.
  				\item Ta tìm khoảng thời gian $\Delta t$ để con lắc đi được quãng đường 1 cm.
  				\item Ta biểu diễn vị trí ban đầu $M_0$ trên đường tròn lượng giác như hình vẽ.
  				\item Để đi được quãng đường 1 cm thì vật quay từ $M_0$ đến M' $\Rightarrow \Delta \varphi = \dfrac{\pi}{6} \Rightarrow \Delta t = \dfrac{T}{12}$.
  				\item Thời gian con lắc đi được 5 cm từ vị trí cân bằng là $t = \dfrac{T}{2} + \dfrac{T}{12} = \dfrac{7T}{12} = \dfrac{7}{4}\ s$.
  		\end{itemize}}{\begin{tikzpicture}[scale=1,thick,>=stealth']
  				\tkzInit[ymin=-3,ymax=3,xmin=-3,xmax=3]\tkzClip
  				\tkzDefPoints{-2.5/0/m, 2.5/0/n, 0/0/o, 0/2.5/p, 0/-2.5/q}
  				\tkzDefShiftPoint[o](-90:2){0}
  				\tkzDefShiftPoint[o](-60:2){m'}
  				\tkzDefPointBy[projection=onto m--n](m') \tkzGetPoint{h}
  				\draw(o)circle(2cm);
  				\tkzDrawSegments[dashed](m',h)
  				\tkzDrawSegments(p,q)
  				\tkzDrawSegments[->](m,n)
  				\tkzDrawSegments[blue,->](o,0 o,m')
  				\tkzDrawPoints[size=3](o,h,0,m')
  				\tkzMarkAngles[size=0.7cm,arrows=->](0,o,m')
  				\node at (0)[below left]{$M_0$};
  				\node at (m')[below right]{$M'$};
  				\node at (2,0)[above right]{$2$};
  				\node at (h)[above left]{$1$};
  		\end{tikzpicture}}
  	}
  	%\haidong{8}
  \end{vidu} %\dotlineEX{5}
  %\loai{Quãng đường vật đi được từ thời điểm $t_1$ đến $t_2$}
  \begin{vidu}%[3P1-13.1K][Loai1] 
  	Một chất điểm dao động điều hòa với phương trình $x=4\sqrt{2}\cos \left({5\pi t-\dfrac{{3\pi }}{4}}\right)$ (x tính bằng cm; t tính bằng s). Quãng đường chất điểm đi từ thời điểm $t_1=0,1$ s đến thời điểm $t_2=6$ s là
  	\choice
  	{84,4 cm}
  	{333,8 cm}
  	{\True 331,4 cm}
  	{336,1 cm}
  	\loigiai{
  		\begin{itemize}
  			\item Tại $t=0$ ta có $x=4\sqrt{2}\cos \left({5\pi t-\dfrac{{3\pi }}{4}}\right)=-4$ cm; và đang chuyển động theo chiều dương $v>0$.
  			
  		\end{itemize}
  	}
  	%\haidong{8}
  \end{vidu} %\dotlineEX{5}
  %\loai{Cho quãng đường tính vận tốc}
  \begin{comment}%[3P1-13.1G][Loai1]
  	Chất điểm dao động điều hoà có chu kì $\dfrac{\pi}{10}\ s$. Chọn gốc thời gian $t = 0$ khi vật nhỏ qua vị trí cân bằng theo chiều dương. Trong khoảng thời gian $\dfrac{\pi}{20}$  s đầu tiên kể từ $t = 0$, vật đi được quãng đường 4 cm. Vận tốc của vật tại thời điểm $\dfrac{\pi}{15}$  s là
  	\choice
  	{$v = 20\sqrt{3}$ cm/s}
  	{$v = -20\sqrt{3}$ cm/s}
  	{$v = 40$ cm/s}
  	{\True $v = - 20$ cm/s}
  	\loigiai{
  		\begin{itemize}
  			\item Kể từ $t = 0$: $x= O\ (+)$, sau $\dfrac{\pi}{20}s=\dfrac{T}{2}$ vật đi được $2A = 4$ cm $\Rightarrow A = 2$ cm.
  			\item Sau $\dfrac{\pi}{15}s=\dfrac{2T}{3}=\dfrac{T}{4}+\dfrac{T}{4}+\dfrac{T}{6}$, vật thực hiện diễn biến dao động như sau:
  			\begin{center}
  				\begin{tikzpicture}
  					\tkzDefPoints{0/0/o, 5/0/a, 4.325/0/b, 3.525/0/c, 2.5/0/d,  -5/0/h, -4.325/0/g, -3.525/0/f, -2.5/0/e, -6/0/x', 6/0/x}
  					\tkzDefShiftPoint[a](-90:0.5cm){a1}
  					\tkzDefShiftPoint[a](-90:1.2cm){a2}
  					\tkzDefShiftPoint[o](-90:0.5cm){o1}
  					\tkzDefShiftPoint[g](-90:1.2cm){g2}
  					\tkzDrawSegments[blue,->,very thick](x',x)
  					\tkzDrawPoints[size=5,fill=black](o,a,b,c,d,e,f,g,h)
  					\tkzLabelPoint[above](o){\tiny $O$}
  					\tkzLabelPoint[above](a){\tiny$A$}
  					\tkzLabelPoint[above](b){\tiny$\dfrac{A\sqrt{3}}{2}$}
  					\tkzLabelPoint[above](c){\tiny$\dfrac{A\sqrt{2}}{2}$}
  					\tkzLabelPoint[above](d){\tiny$\dfrac{A}{2}$}
  					\tkzLabelPoint[above](h){\tiny$-A$}
  					\tkzLabelPoint[above](g){\tiny $-\dfrac{A\sqrt{3}}{2}$}
  					\tkzLabelPoint[above](f){\tiny$-\dfrac{A\sqrt{2}}{2}$}
  					\tkzLabelPoint[above](e){\tiny$-\dfrac{A}{2}$}
  					\tkzLabelPoint[above](x){\tiny$x$}
  					\tkzDrawSegments[->,very thick,red](o1,a1 a2,g2)
  					\tkzDrawSegments[very thick,red](a1,a2)
  				\end{tikzpicture}
  			\end{center}
  			\item Vậy tại thời điểm $\dfrac{\pi}{15}$ s, vật có $x =-\dfrac{A\sqrt{3}}{2}(-) \Rightarrow v=-\dfrac{{{v}_{\max }}}{2}=-20$ cm/s.
  		\end{itemize}
  	}
  	%\haidong{8}
  \end{comment}%\dotlineEX{5}
  %\loai{Tìm quãng đường vật đi được trong giây thứ n}
  \begin{vidu}%[3P1-13.1K][Loai1]
  	Cho một chất điểm dao động điều hòa với phương trình li độ $x = 8\cos(\dfrac{\pi}{2}t)$ cm. Quãng đường chất điểm đi được trong giây thứ 105 là 
  	\choice
  	{$4\sqrt{2}$ cm}
  	{$8\sqrt{2}$ cm}
  	{\True 8 cm}
  	{$8\sqrt{3}$ cm}
  	\loigiai{
  		\immini{\begin{itemize}
  				\item 
  				Ta biểu diễn vị trí pha ban đầu của vật tại vị trí điểm $M_0$ trên đường tròn lượng giác như hình vẽ.
  				\item Ta có: $T = \dfrac{2\pi}{\dfrac{\pi}{2}} = 4+s\Rightarrow t = 105\ s = 104+1 = 26.4+1$.
  				\item Sau 26T, chất điểm quay được 26 vòng và quay về đúng vị trí $M_0 \Rightarrow$ Sau 104 s đầu tiên, vật ở vị trí $M_0$
  				$\Rightarrow$ 1s tiếp theo $= \dfrac{T}{4}$, vật quay thêm góc $\dfrac{\pi}{2}$ đến vị trí $M_1 \Rightarrow$ Quãng đường chất điểm đi được trong giây thứ 105 là $s = 8$ s.
  		\end{itemize}}{\begin{tikzpicture}[scale=1,thick,>=stealth']
  				\tkzInit[ymin=-3,ymax=3,xmin=-3,xmax=3]\tkzClip
  				\tkzDefPoints{-2.5/0/m, 2.5/0/n, 0/0/o, 0/2.5/p, 0/-2.5/q}
  				\tkzDefShiftPoint[o](0:2){0}
  				\tkzDefShiftPoint[o](90:2){m'}
  				\draw(o)circle(2cm);
  				\tkzDrawSegments(p,q)
  				\tkzDrawSegments[->](m,n)
  				\tkzDrawSegments[blue,->](o,0 o,m')
  				\tkzDrawPoints[size=3](o,0,m')
  				\tkzMarkAngles[size=0.7cm,arrows=->](0,o,m')
  				\node at (0)[above right]{$M_0$};
  				\node at (m')[above right]{$M'$};
  				\node at (2,0)[below right]{$8$};
  		\end{tikzpicture}}
  	}
  	%\haidong{7}
  \end{vidu} %\dotlineEX{5}
  
  %\loai{Bài toán tổng hợp}
  \begin{comment} %[3P1-13.2G][Loai1] 
  	\nguon{QG - 2018} Một vật dao động điều hòa quanh vị trí cân bằng O. Tại thời điểm $t_1$, vật đi qua vị trí cân bằng. Trong khoảng thời gian từ thời điểm $t_1$ đến thời điểm $t_2 = t_1 + \dfrac{1}{6}\ s$, vật không đổi chiều chuyển động và tốc độ của vật giảm còn một nửa. Trong khoảng thời gian từ thời điểm $t_2$ đến thời điểm $t_3 = t_2 + \dfrac{1}{6}\ s$, vật đi được quãng đường 6 cm. Tốc độ cực đại của vật trong quá trình dao động là
  	\choice
  	{37,7 m/s}
  	{0,38 m/s}
  	{\True 1,41 m/s}
  	{224 m/s}
  	\loigiai{
  		\begin{itemize}
  			\item Thời gian tốc độ giảm 1 nửa: $\dfrac{1}{6}=\dfrac{T}{6}\Rightarrow T=1\ s\Rightarrow \omega =2\pi\ rad/s;\  v=\dfrac{1}{2}v_{\max}\Rightarrow x=\pm \dfrac{A\sqrt{3}}{2}$.
  			\item vẽ vòng tròn lượng giác ta xác định được quãng đường đi từ $t_2$ đến $t_3$: 
  			\begin{align*}
  				S= A-\dfrac{A\sqrt{3}}{2}=3\ cm\Rightarrow A=12+6\sqrt{3}\ cm\Rightarrow v_{\max}=\left(12+6\sqrt{3}\right)2\pi =140,695\ cm/s.
  			\end{align*}
  			
  		\end{itemize}
  	}
  	%\haidong{8}
  \end{vidu} %\dotlineEX{9}
  
  %\loai{Tốc độ trung bình trong khoảng thời gian từ $t_1$ đến $t_2$}
  \begin{vidu}%[3P1-13.3K][Loai1] 
  	Cho một chất điểm dao động điều hòa theo phương trình $x = 6\cos(0,5\pi t + \dfrac{\pi}{6})$, x tính bằng cm và t tính bằng s. Tính từ thời điểm ban đầu, tốc độ trung bình của vật trong 7 s chuyển động là
  	\choice
  	{5,36 cm/s}
  	{3,93 cm/s}
  	{5,14 cm/s}
  	{\True 6,31 cm/s}
  	\loigiai{
  		\begin{itemize}
  			\item $T = \dfrac{2\pi}{0,5\pi} = 4\ s$.
  			\item Ta có: $t = 7\ s = 1.4 + 3\Rightarrow \alpha = 1.2\pi+ \dfrac{3\pi}{2}$.
  			\item Tại thời điểm $t = 0$, vật ở vị trí $M_0$, sau 1 chu kì, chất điểm quay được đúng 1 vòng đến vị trí $M_0$.
  		\end{itemize}
  		\immini{
  			\begin{itemize}
  				\item Trong khoảng thời gian $\Delta t = 3\ s$ tiếp, chất điểm quay được góc $\dfrac{3\pi}{2}$ đến vị trí M'.
  				$\Rightarrow$ quãng đường đi được trong khoảng thời gian $\Delta t$ là $s = 3\sqrt{3}+2.6+3 = 15+3\sqrt{3}$.\\ 
  				$\Rightarrow$ Quãng đường chất điểm đi được trong 7 s là $4.6+15+3\sqrt{3} = 39+3\sqrt{3}$ cm.\\
  				$\Rightarrow$ Vận tốc trung bình trong 7s là $\dfrac{39+3\sqrt{3}}{7} = 6,31$ cm/s.
  				
  			\end{itemize}
  		}{\begin{tikzpicture}[scale=1,thick,>=stealth']
  				\tkzInit[ymin=-3,ymax=3,xmin=-3,xmax=3]\tkzClip
  				\tkzDefPoints{-2.5/0/m, 2.5/0/n, 0/0/o, 0/2.5/p, 0/-2.5/q}
  				\tkzDefShiftPoint[o](30:2){0}
  				\tkzDefShiftPoint[o](-60:2){m'}
  				\tkzDefPointBy[projection=onto m--n](0) \tkzGetPoint{h}
  				\tkzDefPointBy[projection=onto m--n](m') \tkzGetPoint{h'}
  				\draw(o)circle(2cm);
  				\tkzDrawSegments[dashed](m',h' 0,h)
  				\tkzDrawSegments(p,q)
  				\tkzDrawSegments[->](m,n)
  				\tkzDrawSegments[blue,->](o,0 o,m')
  				\tkzDrawPoints[size=3](o,h,0,m',h')
  				\tkzMarkAngles[size=0.7cm,arrows=->](0,o,m')
  				\node at (0)[above right]{$M_0$};
  				\node at (m')[below right]{$M'$};
  				\node at (2,0)[below right]{$6$};
  				\node at (h')[below]{$3$};
  				\node at (h)[below]{$3\sqrt{3}$};
  		\end{tikzpicture}}
  	}
  	%\haidong{6}
  \end{comment} %\dotlineEX{4}
  
  %\loai{Tốc độ trung bình trong khoảng thời gian vật đi từ $x_1$ đến $x_2$}
  \begin{vidu}%[3P1-13.3K][Loai1] 
  	Một chất điểm dao động trên trục Ox quanh vị trí cân bằng chính là gốc tọa độ O. Quỹ đạo của vật là đoạn thẳng dài 4 cm. Quãng thời gian ngắn nhất để vật đi từ li độ $x_1 = 1$ cm đến $x_2 = 2$ cm là 0,5 s. Tốc độ trung bình của vật trong một chu kì dao động là
  	\choice
  	{2 cm/s}
  	{4 cm/s}
  	{$\dfrac{4}{3}$ cm/s}
  	{\True $\dfrac{8}{3}$ cm/s}
  	\loigiai{
  		\immini{
  			Quãng thời gian ngắn nhất để vật đi từ li độ $x_1 = 1$ cm đến $x_2 = 2$ cm ứng với vật đi từ $M_0$ đến $M_1\Rightarrow \dfrac{T}{6} = 0,5 \Rightarrow t = 3\ s\Rightarrow v_{tb} = \dfrac{4.2}{3} = \dfrac{8}{3}\ m/s$.
  		}{\begin{tikzpicture}[scale=1,thick,>=stealth']
  				\tkzInit[ymin=-3,ymax=3,xmin=-3,xmax=3]\tkzClip
  				\tkzDefPoints{-2.5/0/m, 2.5/0/n, 0/0/o, 0/2.5/p, 0/-2.5/q}
  				\tkzDefShiftPoint[o](-60:2){0}
  				\tkzDefShiftPoint[o](0:2){m'}
  				\tkzDefPointBy[projection=onto m--n](0) \tkzGetPoint{h}
  				\draw(o)circle(2cm);
  				\tkzDrawSegments[dashed](h,0)
  				\tkzDrawSegments(p,q)
  				\tkzDrawSegments[->](m,n)
  				\tkzDrawSegments[blue,->](o,0 o,m')
  				\tkzDrawPoints[size=3](o,h,0,m')
  				\tkzMarkAngles[size=0.7cm,arrows=->](0,o,m')
  				\node at (0)[below right]{$M_0$};
  				\node at (m')[above right]{$M_1$};
  				\node at (2,0)[below right]{$2$};
  				\node at (h)[below right]{$1$};
  		\end{tikzpicture}}
  	}
  	%\haidong{5}
  \end{vidu} %\dotlineEX{11}
  %\loai{Tốc độ trung bình trong khoảng thời gian vật đi được quãng đường s}
  \begin{vidu}%[3P1-13.3K][Loai1]
  	Một chất điểm dao động với phương trình  $x=10\cos \left( 2\pi t-\dfrac{2\pi}{3} \right)$ cm  (t tính bằng s). Tốc độ trung bình của chất điểm khi nó đi được quãng đường 70 cm đầu tiên (kể từ $t = 0$) là
  	\choice
  	{50 cm/s}
  	{40 cm/s}
  	{35 cm/s}
  	{\True 42 cm/s}
  	\loigiai{
  		
  	}
  	%\haidong{4}
  \end{vidu} %\dotlineEX{11}  
  %\loai{Cho khoảng thời gian thỏa mãn tốc độ và tốc độ trung bình. Tìm các đại lượng}
  \begin{vidu}%[3P1-13.3K][Loai1]
  	Một vật nhỏ dao động điều hòa. Khoảng thời gian giữa hai lần liên tiếp vật có tốc độ bằng không là 1 s, đồng thời tốc độ trung bình trong khoảng thời gian này là 10 cm/s. Khi qua vị trí cân bằng, tốc độ của vật nhỏ là
  	\choice
  	{\True 15,7 cm/s}
  	{31,4 cm/s}
  	{20 cm/s}
  	{10 cm/s}
  	\loigiai{
  		\begin{itemize}
  			\item Vật có tốc độ bằng 0 tại 2 biên $\Rightarrow $ hai lần liên tiếp vật có tốc độ bằng 0 khi vật đi từ biên này sang biên kia
  			$\Rightarrow $ quãng đường vật đi là $S = 2A$, thời gian đi là $\Delta t=\dfrac{T}{2}=1\ s\Rightarrow   T=2\ s$.
  			\item Mà $v_{tb}=\dfrac{S}{\Delta t}\Rightarrow   S= 10\ cm = 2A \Rightarrow  A = 5$ cm.
  			\item Vậy khi qua vị trí cân bằng, tốc độ vật nhỏ là $v_{\max}=\omega A = 15,7$ cm.
  			
  		\end{itemize}
  	}
  	%\haidong{6}
  \end{vidu} %\dotlineEX{11}
  %\loai{Tốc độ trung bình giữa hai thời điểm có li độ và gia tốc cho trước}
  \begin{comment}%[3P1-13.4K][Loai1] 
  	Một vật nhỏ dao động điều hòa theo một quỹ đạo thẳng dài 20 cm với chu kì 1 s. Từ thời điểm vật qua vị trí có li độ $5\sqrt{3}$ cm theo chiều dương đến khi gia tốc của vật đạt giá trị cực tiểu lần thứ hai, vật có tốc độ trung bình xấp xỉ bằng
  	\choice
  	{58,7 cm/s}
  	{\True 38,2 cm/s}
  	{31,3 cm/s}
  	{41,3 cm/s}
  	\loigiai{
  		\immini{\begin{itemize}
  				\item 
  				Ta có: $2A = 20 \Rightarrow  A = 10\ cm,\ T = 1$ s.
  				\item Ta biểu diễn vị trí có li độ $5\sqrt{3}$ cm theo chiều dương bằng điểm $M_0$ trên hình vẽ. 
  				\item Gia tốc của vật có giá trị cực tiểu tại biên dương (vị trí điểm M' trên đường tròn).
  				\item Từ hình vẽ, ta suy ra quãng đường vật đi từ thời điểm vật qua vị trí có li độ $5\sqrt{3}$ cm theo chiều dương đến khi gia tốc của vật đạt giá trị cực tiểu lần thứ hai là $s = 10 - 5\sqrt{3}\ cm + 10.4 = 50 - 5\sqrt{3}\ cm$.
  		\end{itemize}}{\begin{tikzpicture}[scale=1,thick,>=stealth']
  				\tkzInit[ymin=-3,ymax=3,xmin=-3,xmax=3]\tkzClip
  				\tkzDefPoints{-2.5/0/m, 2.5/0/n, 0/0/o, 0/2.5/p, 0/-2.5/q}
  				\tkzDefShiftPoint[o](-30:2){0}
  				\tkzDefShiftPoint[o](0:2){m'}
  				\tkzDefPointBy[projection=onto m--n](0) \tkzGetPoint{h}
  				\draw(o)circle(2cm);
  				\tkzDrawSegments[dashed](0,h)
  				\tkzDrawSegments(p,q)
  				\tkzDrawSegments[->](m,n)
  				\tkzDrawSegments[blue,->](o,0 o,m')
  				\tkzDrawPoints[size=3](o,h,0,m')
  				\tkzMarkAngles[size=0.7cm,arrows=->](0,o,m')
  				\node at (0)[below right]{$M_0$};
  				\node at (m')[below right]{$M'$};
  				\node at (2,0)[above right]{$10$};
  				\node at (h)[above left]{$5\sqrt{3}$};
  		\end{tikzpicture}}
  		\begin{itemize}
  			\item Thời gian để vật qua vị trí có li độ $5\sqrt{3}$ cm theo chiều dương đến khi gia tốc của vật đạt giá trị cực tiểu lần thứ hai là $\dfrac{T}{12} + T = 13\dfrac{T}{12} = \dfrac{13}{12}$ s $\Rightarrow $ tốc độ trung bình của vật là $v_{tb} = \dfrac{50-5\sqrt{3}}{\dfrac{13}{12}} = 38,2$ cm/s.
  		\end{itemize}
  	}
  	%\haidong{8}
  \end{vidu} %\dotlineEX{12}
  %\loai{Tốc độ trung bình giữa hai thời điểm có vận tốc và gia tốc cho trước}
  \begin{vidu}%[3P1-13.4K][Loai1]
  	Một vật nhỏ dao động điều hòa theo một quỹ đạo thẳng dài 10 cm với chu kì 2 s. Cho $\pi ^2 = 10$. Từ thời điểm vật qua vị trí có gia tốc $-25\ cm/s^2$ theo chiều âm đến khi vận tốc của vật đạt giá trị cực đại lần thứ 5, vật có tốc độ trung bình là
  	\choice
  	{12,33 cm/s}
  	{12,73 cm/s}
  	{\True 10,09 cm/s}
  	{11,32 cm/s}
  	\loigiai{
  		\begin{itemize}
  			\item Vật theo chiều âm và có $a=-\omega^2x=-25c\ m/s^2\Leftrightarrow x=\dfrac{A}{2}\ (-)$.
  			\item Vận tốc đạt giá trị cực đại $\Leftrightarrow x = 0\ (+)$.
  			\item Một chu kì, vật qua $x = 0\ (+)$ 1 lần. Tách $5 = 4 + 1$.
  			\item Sau 4T, vật đi được 16A và qua $x = 0\ (+)$ 4 lần và trở lại trạng thái $x=\dfrac{A}{2}\ (-)$. Vật đi thêm 1 lần nữa mất $\dfrac{T}{3}+\dfrac{T}{4}$ và đi được thêm 2,5A.
  			\begin{center}
  				\begin{tikzpicture}
  					\tkzDefPoints{0/0/o, 5/0/a, 4.325/0/b, 3.525/0/c, 2.5/0/d,  -5/0/h, -4.325/0/g, -3.525/0/f, -2.5/0/e, -6/0/x', 6/0/x}
  					\tkzDefShiftPoint[d](-90:0.5cm){d1}
  					\tkzDefShiftPoint[h](-90:1cm){h2}
  					\tkzDefShiftPoint[h](-90:0.5cm){h1}
  					\tkzDefShiftPoint[o](-90:1cm){o2}
  					\tkzDrawSegments[blue,->,very thick](x',x)
  					\tkzDrawPoints[size=5,fill=black](o,a,b,c,d,e,f,g,h)
  					\tkzLabelPoint[above](o){\tiny $O$}
  					\tkzLabelPoint[above](a){\tiny$A$}
  					\tkzLabelPoint[above](b){\tiny$\dfrac{A\sqrt{3}}{2}$}
  					\tkzLabelPoint[above](c){\tiny$\dfrac{A\sqrt{2}}{2}$}
  					\tkzLabelPoint[above](d){\tiny$\dfrac{A}{2}$}
  					\tkzLabelPoint[above](h){\tiny$-A$}
  					\tkzLabelPoint[above](g){\tiny $-\dfrac{A\sqrt{3}}{2}$}
  					\tkzLabelPoint[above](f){\tiny$-\dfrac{A\sqrt{2}}{2}$}
  					\tkzLabelPoint[above](e){\tiny$-\dfrac{A}{2}$}
  					\tkzLabelPoint[above](x){\tiny$x$}
  					\tkzDrawSegments[->,very thick,red](d1,h1 h2,o2)
  					\tkzDrawSegments[very thick,red](h1,h2)
  				\end{tikzpicture}
  			\end{center}
  			\item Vậy tốc độ trung bình vật trong quá trình này là: $v_{tb}=\dfrac{S}{\Delta t}=\dfrac{16A+2,5A}{4T+\dfrac{T}{3}+\dfrac{T}{4}}=10,09$ cm/s.
  		\end{itemize}
  	}
  	%\haidong{8}
  \end{comment} %\dotlineEX{11}
  
  %\loai{Tốc độ trung bình từ thời điểm ban đầu tới thời điểm $W_{\text{đ}}=nW_t$}
  \begin{comment}%[3P1-11.2K][Loai1]%Bài 13
  	Một vật dao động điều hoà với phương trình vận tốc $ v=10\pi \cos  \left( \pi t+\dfrac{\pi}{3} \right) $  cm/s. Tốc độ trung bình của vật trên quãng đường từ thời điểm ban đầu tới thời điểm vật có động năng bằng 3 lần thế năng lần thứ 3 là
  	\choice
  	{13,33 cm/s}
  	{\True 17,56 cm/s}
  	{15 cm/s}
  	{20 cm/s}
  	\loigiai{
  		\begin{itemize}
  			\item Trạng thái: $W_{\text{đ}}=3{W_t}\Leftrightarrow x=\pm \dfrac{A}{2}.$ 
  			\item Tại $t = 0:\ \varphi _v = \dfrac{\pi}{3} \Rightarrow  {\varphi _x} = - \dfrac{\pi}{6} \Rightarrow  x = \dfrac{A\sqrt 3}{2} \oplus.$
  			\item Vậy: $ {{v}_{tb}}=\dfrac{2,5A+A-\dfrac{A\sqrt{3}}{2}}{\dfrac{T}{12}+\dfrac{T}{2}+\dfrac{T}{6}}=  17,56\ cm/s$.
  			
  		\end{itemize}
  	}
  	%\haidong{8}
  \end{comment} %\dotlineEX{13}
  
  
  
  
  